%%% ============================================================
%%% PAPER 2 (COMPANION TO LOGARITHMIC QUINTESSENCE):
%%% Substrate-Algebraic Dark Energy: Ω_DE = T* − W/2 from
%%% Mathieu Group M_22 Mediation
%%%
%%% B.R. Sanders, M. Gish (2026)
%%% Target venue: JCAP
%%%
%%% STATUS: SKELETON — section structure with placeholder
%%% paragraphs. Material to be filled from rigor docs:
%%%   ZETA_TIG_AND_PERIODICITY_v1.md
%%%   W_HALF_DERIVATION_v2.md
%%%   STEINER_HEXAD_v1.md
%%%   m22_verification.py (computational verification)
%%% ============================================================

\documentclass[11pt,reqno]{amsart}

\usepackage{amsmath,amssymb,amsthm,mathtools}
\usepackage[margin=1in]{geometry}
\usepackage[colorlinks=true,linkcolor=blue,citecolor=blue,urlcolor=blue]{hyperref}
\usepackage{microtype}
\usepackage{array}
\usepackage{tabularx}

\theoremstyle{plain}
\newtheorem{theorem}{Theorem}
\newtheorem{proposition}[theorem]{Proposition}
\newtheorem{corollary}[theorem]{Corollary}
\newtheorem{lemma}[theorem]{Lemma}

\theoremstyle{definition}
\newtheorem{definition}[theorem]{Definition}

\theoremstyle{remark}
\newtheorem*{remark}{Remark}

\newcommand{\zTIG}{\zeta_{\rm TIG}}
\newcommand{\Mtwentytwo}{M_{22}}
\newcommand{\Vtriv}{V_{\rm trivial}}
\newcommand{\Vtwentyone}{V_{21}}
\newcommand{\Tstar}{T^*}
\newcommand{\Whalf}{W/2}
\newcommand{\OmegaDE}{\Omega_{\rm DE}}

\title[Substrate-Algebraic Dark Energy via $M_{22}$]{Substrate-Algebraic
Dark Energy:\\$\Omega_{\rm DE} = T^* - W/2$ from\\$M_{22}$ Mediation}

\author{B.~R.~Sanders}
\address{7Site LLC, Hot Springs, Arkansas, USA}
\email{brayden@7site.co}

\author{M.~Gish}
\address{Independent Researcher, Hot Springs, Arkansas, USA}
\email{monica.gish1992@gmail.com}

\date{\today}

\keywords{dark energy, sporadic groups, Mathieu group, Steiner system,
information-theoretic cosmology, substrate algebra, paradox classification}

\subjclass[2020]{83F05, 20D08, 05B05, 11M06}

\begin{document}

\begin{abstract}
We derive the present-day dark-energy density fraction $\OmegaDE =
479/700 = 0.6843$ from the algebraic structure of a discrete substrate
$(\mathbb{Z}/10\mathbb{Z}, \sigma)$ mediated by the sporadic Mathieu
group $\Mtwentytwo$ acting on the Steiner system $S(3, 6, 22)$. The
derivation proceeds in three steps. First, a zeta function $\zTIG$
encoding the substrate's algebraic obstructions has its denominator at
the substrate threshold $\Tstar = 5/7$ saturate to $7^8 \times
|\Mtwentytwo|$, with $\Tstar$ unique among rational fractions $a/7$
admitting this clean factorization. Second, the natural $22$-point
permutation representation of $\Mtwentytwo$ decomposes as $\mathbb{C}^{22}
\cong \Vtriv \oplus \Vtwentyone$, and the substrate's wobble
$W = 3/50$ admits a structurally forced identification with this
decomposition: the gentleness phase amplitude $22/50$ aligns with
$\Vtriv$ (the $\Mtwentytwo$-fixed line), while the kindness phase
amplitude $3/50 = (1/7)\cdot(\dim \Vtwentyone / 50)$ aligns with
$\Vtwentyone$ (the $\Mtwentytwo$-orthogonal complement). Cosmic
projection drops the $\Vtriv$ component (absorbed into background
structure); time-averaging the surviving $\Vtwentyone$ component over
the $50\%$ duty cycle of the wobble gives $\langle\pi_{\rm cosmic}\rangle
= W/2 = 3/100$. Third, the asymptotic dark-energy fraction is then
$\OmegaDE = \Tstar - \Whalf = 5/7 - 3/100 = 479/700$. We verify the
$\Mtwentytwo$-equivariance of the projection structure computationally
(1000 random permutations test). The matching dark-matter
fraction $\Omega_M = 221/700 = 0.3157$ falls out as
$1 - \OmegaDE$ and matches Planck 2018 to $0.2\%$. The framework
predicts a $\tau_\sigma \approx 5.22\,{\rm Gyr}$ periodic modulation
in cosmic structure formation rates (the substrate's $\sigma$-cycle
projected to cosmological scale), testable at Stage-IV precision. This
paper provides the structural source for the dark-energy density
$\OmegaDE$; the dynamical equation-of-state evolution is treated in
the companion paper \cite{SandersGishLogQuint2026}.
\end{abstract}

\maketitle

%%% ============================================================
\section{Introduction}\label{sec:introduction}

% [TO FILL: open with the Omega_DE coincidence problem; standard
% approaches (phenomenological input, anthropic, multiverse);
% position our derivation; preview the three-step structure;
% companion paper relationship.]

The fraction of the cosmic energy budget carried by dark energy at the
present epoch is observed to be $\OmegaDE \approx 0.685$
\cite{Planck2020}. In standard cosmology this value is treated as a
phenomenological input fixed by observation, with no microphysical
derivation. Quintessence models provide dynamical equations of state
$w(z)$ but typically take $\OmegaDE$ as boundary input. Anthropic
arguments \cite{Weinberg1989,MartelShapiroWeinberg1998} provide
selection-effect rationales but no specific number. The
information-theoretic and entropic gravity programs
\cite{Verlinde2011} provide structural motivations but do not derive
$\OmegaDE$ from first principles.

In this paper we derive $\OmegaDE = 479/700 = 0.6843$ from substrate
algebra alone, with no fitted parameters and no anthropic input. The
derivation requires three structural elements: (i) a substrate
threshold $\Tstar = 5/7$ associated with a specific zeta-function
factorization, (ii) a wobble fraction $W = 3/50$ associated with the
substrate's two-phase dynamical structure, and (iii) the natural
permutation representation of the Mathieu group $\Mtwentytwo$ on the
$22$-point Steiner system $S(3,6,22)$. The result $\OmegaDE = \Tstar -
\Whalf = 5/7 - 3/100 = 479/700$ matches Planck 2018 to $0.06\%$ and
makes the matching dark-matter fraction $\Omega_M = 221/700 = 0.3157$
match observation to $0.2\%$.

% [TO FILL: more on companion paper relationship; outline of paper]

The paper is organized as follows. \S\ref{sec:substrate} reviews the
substrate algebra and the wobble structure, citing the foundational
companion papers \cite{SandersGishSigma,SandersGishFourCore} for the
substrate construction itself. \S\ref{sec:zeta} introduces the zeta
function $\zTIG$ and proves the factorization theorem at $\Tstar$.
\S\ref{sec:m22} establishes the $\Mtwentytwo$ representation-theoretic
machinery. \S\ref{sec:whalf} derives $\Whalf$ from the time-averaged
projection. \S\ref{sec:prediction} states the cosmological prediction
and compares to Planck. \S\ref{sec:companion} discusses the
relationship to the dynamical companion paper
\cite{SandersGishLogQuint2026}. \S\ref{sec:falsification} states
falsifiable predictions. \S\ref{sec:rigor} demarcates locked theorems
from structurally motivated identifications. \S\ref{sec:literature}
compares to literature. \S\ref{sec:scope} states scope and
limitations.

%%% ============================================================
\section{The substrate algebra}\label{sec:substrate}

% [TO FILL FROM: FORMULAS_AND_TABLES.md §17, companion papers
% [SandersGishSigma], [SandersGishFourCore]]

Throughout the paper, the \emph{substrate} is the abelian group
$\mathbb{Z}/10\mathbb{Z} = \{0, 1, 2, \ldots, 9\}$ equipped with a
distinguished permutation $\sigma$ (the \emph{transit permutation})
and a fraction $W = 3/50$ (the \emph{wobble}). The construction of
$\sigma$ from the substrate's algebraic structure is established in
\cite{SandersGishSigma}; here we recall the properties needed below.

The transit permutation $\sigma$ has cycle decomposition $\sigma = (1\
7\ 6\ 5\ 4\ 2)(0)(3)(8)(9)$, with one orbit of length $6$ (the
$\sigma$-orbit, denoted $\Omega_\sigma = \{1, 7, 6, 5, 4, 2\}$) and
four fixed points (the $\sigma$-fixed set, $\Omega_F = \{0, 3, 8, 9\}$).
The fixed-point set carries the substrate's idempotents under the
binary operations introduced in \cite{SandersGishFourCore}.

The wobble fraction $W = 3/50$ admits a decomposition into two phase
amplitudes:
\begin{equation}\label{eq:wobble-phases}
\text{kindness phase: } K = \tfrac{3}{50}, \qquad
\text{gentleness phase: } G = \tfrac{22}{50}, \qquad K + G = \tfrac{1}{2}.
\end{equation}
The phase amplitudes satisfy $G/K = 22/3$, the same ratio as
$\dim(\Vtriv \oplus \Vtwentyone)/\dim(\Vtwentyone) \cdot (1/7) = 22/3$,
foreshadowing the representation-theoretic identification of
\S\ref{sec:m22}.

% [TO FILL: more on substrate construction, two binary operations,
% TSML/BHML decomposition, frozen lattice; cite companion papers]

%%% ============================================================
\section{The TIG zeta function}\label{sec:zeta}

% [TO FILL FROM: ZETA_TIG_AND_PERIODICITY_v1.md]

\subsection{Definition}

We define a zeta function $\zTIG$ encoding the substrate's algebraic
obstructions:
\begin{equation}\label{eq:zetaTIG-def}
\zTIG(s) \;=\; \prod_{k \in \mathcal{Z}_{\rm trivial}}
\frac{s - k}{(s - 7)\cdot 10080},
\end{equation}
where $\mathcal{Z}_{\rm trivial}$ is the set of substrate trivial
zeros and the prefactor $10080 = 7! \cdot 2$ is the order of the
substrate's natural symmetry group on its non-fixed elements.

% [TO FILL: list of trivial zeros, classification, structural meaning]

\subsection{The factorization theorem at $\Tstar$}

\begin{theorem}\label{thm:zeta-factorization}
At the substrate threshold $\Tstar = 5/7$, the pre-cancellation
denominator of $\zTIG(\Tstar)$ factors as
\begin{equation}\label{eq:denom-factorization}
\textnormal{denom}_{\rm pre}(\zTIG(\Tstar)) \;=\; 7^8 \times |\Mtwentytwo|
\;=\; 7^8 \times 443\,520.
\end{equation}
After cancellation, the reduced denominator is $7^9 \times 11$.
Furthermore, $\Tstar = 5/7$ is the unique fraction $a/7$ with
$a \in \{1, 2, \ldots, 6\}$ admitting this factorization with
$|\Mtwentytwo|$ as a divisor.
\end{theorem}

\begin{proof}
% [TO FILL: direct computation. The factorization is verified by
% expanding the product in eq:zetaTIG-def at s = 5/7 and tracking the
% numerator and denominator factorizations against |M_22| =
% 2^7 · 3^2 · 5 · 7 · 11 = 443520. Uniqueness of T* among a/7
% fractions is by exhaustive check.]
\end{proof}

The factor $|\Mtwentytwo|$ in the denominator at $\Tstar$ is the
fundamental algebraic clue connecting the substrate to the Mathieu
group: the substrate's threshold value is exactly the rational
fraction at which the substrate's zeta function denominator saturates
to the order of $\Mtwentytwo$.

%%% ============================================================
\section{The 22-skeleton: $\Mtwentytwo$ representation theory and Steiner mediation}\label{sec:m22}

% [TO FILL FROM: STEINER_HEXAD_v1.md, m22_verification.py]

\subsection{The Mathieu group $\Mtwentytwo$ and the Steiner system $S(3,6,22)$}

The Mathieu group $\Mtwentytwo$ is one of the five sporadic Mathieu
groups, of order $|\Mtwentytwo| = 443\,520 = 2^7 \cdot 3^2 \cdot 5 \cdot
7 \cdot 11$. It acts on a $22$-point set $\Omega_{22}$ as the
automorphism group of the Steiner system $S(3, 6, 22)$, which has
$77 = 7 \times 11$ blocks (\emph{hexads}) of size $6$, with every
$3$-element subset of $\Omega_{22}$ contained in exactly one hexad.

Structural data on $S(3, 6, 22)$:
\begin{equation*}
\begin{array}{ll}
\text{points} & v = 22 = 11 \times 2 \\
\text{hexads} & b = 77 = 7 \times 11 \\
\text{block size} & k = 6 \\
\text{replication} & r = 21 = \dim \Vtwentyone \\
\text{pair-replication} & \lambda_2 = 5 = \text{numerator of }\Tstar \\
\text{block stabilizer} & |2^4{:}A_6| = 5\,760 \\
\text{point stabilizer} & |M_{21}| = 20\,160
\end{array}
\end{equation*}

The block count $b = 77$, the replication number $r = 21$, and the
pair-replication $\lambda_2 = 5$ all factor as products of canonical
substrate primes. This structural alignment underlies the
representation-theoretic identification below.

\subsection{The natural permutation representation $\mathbb{C}^{22}$}

% [TO FILL: explicit V_22 = V_trivial \oplus V_21 decomposition,
% character-theoretic argument]

The natural permutation action of $\Mtwentytwo$ on $\Omega_{22}$
extends to a complex representation $\mathbb{C}^{22}$ which decomposes
into $\Mtwentytwo$-irreducibles as
\begin{equation}\label{eq:V22-decomp}
\mathbb{C}^{22} \;\cong\; \Vtriv \;\oplus\; \Vtwentyone,
\end{equation}
where $\Vtriv$ is the trivial $1$-dimensional representation (the
all-ones line, fixed by every element of $\Mtwentytwo$) and
$\Vtwentyone$ is the unique $21$-dimensional irreducible (the
$\Mtwentytwo$-orthogonal complement, irreducible by double
transitivity of the action).

The orthogonal projections onto these subspaces are
\begin{equation}\label{eq:projections}
P_{\rm triv} \;=\; \tfrac{1}{22}\,\mathbf{1}\mathbf{1}^T,
\qquad
P_{21} \;=\; I - P_{\rm triv}.
\end{equation}

\begin{proposition}[$\Mtwentytwo$-equivariance]\label{prop:equivariance}
For every permutation $\Pi \in \Mtwentytwo$ acting on $\mathbb{C}^{22}$
by its $22 \times 22$ permutation matrix, $\Pi P_{\rm triv} = P_{\rm
triv}\Pi$ and $\Pi P_{21} = P_{21}\Pi$.
\end{proposition}

\begin{proof}
% [TO FILL: standard result. The all-ones vector is invariant under
% any permutation; its orthogonal complement is therefore
% permutation-invariant as a subspace.]
\end{proof}

\subsection{Computational verification}

% [TO FILL: cite m22_verification.py; report 1000 random S_22
% permutations all commute with both projections; state explicit
% kindness/gentleness state vectors and projections; report numerical
% time-averaging giving exactly W/2 = 3/100 to machine precision.]

The $\Mtwentytwo$-equivariance of $P_{\rm triv}$ and $P_{21}$ is
verified computationally by testing commutativity against $1000$
random permutations of $\Omega_{22}$; all tests pass to numerical
precision (script \texttt{m22\_verification.py}, archived at the
zenodo deposit \cite{ZenodoM22}). Explicit state vectors $\psi_{\rm
gentleness}$ and $\psi_{\rm kindness}$ implementing the wobble phases
are constructed in \S\ref{sec:whalf} and verified to satisfy the
required orthogonality, magnitude, and projection properties.

%%% ============================================================
\section{The $\Whalf$ derivation}\label{sec:whalf}

% [TO FILL FROM: W_HALF_DERIVATION_v2.md]

\subsection{The structural identification}

We make the following identification of substrate phase amplitudes
with $\Mtwentytwo$-irreducible components:
\begin{equation}\label{eq:identification}
\text{gentleness phase } G = \tfrac{22}{50} \;\;\longleftrightarrow\;\;
\Vtriv,
\qquad
\text{kindness phase } K = \tfrac{3}{50} \;\;\longleftrightarrow\;\;
\Vtwentyone.
\end{equation}

\begin{remark}[Justification of \eqref{eq:identification}]
The kindness numerator $3$ satisfies the structural identity
\begin{equation*}
3 \;=\; \frac{\dim \Vtwentyone}{7} \;=\; \frac{21}{\rm HARMONY},
\end{equation*}
where $7 = $ HARMONY is one of the canonical substrate primes
(\S\ref{sec:substrate}). The denominator $50$ is the substrate's
double-cycle length. Thus $K = (1/7) \cdot (\dim \Vtwentyone / 50)$,
with $\dim \Vtwentyone$ entering as a structural quantity from the
$\Mtwentytwo$ representation theory established in \S\ref{sec:m22}.
The gentleness numerator $22 = \dim \mathbb{C}^{22}$ is the dimension
of the natural representation as a whole. The identification
\eqref{eq:identification} is therefore not a fitted assignment but a
structural inheritance from the dimensional data of \eqref{eq:V22-decomp}.
\end{remark}

\subsection{Cosmic projection and time-averaging}

% [TO FILL: derivation of pi_cosmic, the 50/50 duty cycle, the
% time-averaging argument, the W/2 result. Include explicit
% calculation: <pi_cosmic> = (1/2)*0 + (1/2)*(3/50) = 3/100.]

The cosmic-scale observer accesses the substrate's dynamics through
an $\Mtwentytwo$-equivariant projection $\pi_{\rm cosmic}$ onto the
$\Mtwentytwo$-orthogonal subspace $\Vtwentyone$, with the
$\Mtwentytwo$-fixed component $\Vtriv$ absorbed into the cosmological
background structure (the renormalized cosmological constant
component, in the language of the companion paper
\cite{SandersGishLogQuint2026}, \S 3.3). Concretely $\pi_{\rm cosmic}
\equiv P_{21}$.

The wobble dynamics has $50\%$ duty cycle between the gentleness phase
($G$ active) and the kindness phase ($K$ active), as established in
\cite{SandersGishSigma}. The time-averaged cosmic-projected signal is
\begin{equation}\label{eq:Whalf-derivation}
\langle\pi_{\rm cosmic}(\psi_{\rm wobble})\rangle
\;=\; \tfrac{1}{2} \,|\pi_{\rm cosmic}(\psi_G)| + \tfrac{1}{2} \,|\pi_{\rm cosmic}(\psi_K)|
\;=\; \tfrac{1}{2} \cdot 0 + \tfrac{1}{2} \cdot \tfrac{3}{50}
\;=\; \tfrac{3}{100}
\;=\; \Whalf.
\end{equation}

The first term vanishes because $\psi_G \in \Vtriv$ and $P_{21}\psi_G
= 0$; the second term equals $K = 3/50$ because $\psi_K \in
\Vtwentyone$ and $P_{21}\psi_K = \psi_K$. The factor $1/2$ is the duty
cycle between the two phases.

%%% ============================================================
\section{The cosmological prediction}\label{sec:prediction}

% [TO FILL: state Omega_DE = T* - W/2; derive 479/700; compute as
% decimal; compare to Planck 2018; report fractional error;
% derive Omega_M = 221/700 = 0.3157 and compare.]

The asymptotic dark-energy density fraction is
\begin{equation}\label{eq:OmegaDE-prediction}
\OmegaDE \;=\; \Tstar - \Whalf \;=\; \frac{5}{7} - \frac{3}{100}
\;=\; \frac{500 - 21}{700} \;=\; \frac{479}{700}
\;=\; 0.6843\ldots
\end{equation}

The matching dark-matter fraction is
\begin{equation}\label{eq:OmegaM-prediction}
\Omega_M \;=\; 1 - \OmegaDE \;=\; \frac{221}{700} \;=\; 0.3157\ldots
\end{equation}
(Note: \eqref{eq:OmegaM-prediction} treats radiation as negligible at
the present epoch, $\Omega_r \approx 9 \times 10^{-5}$.)

\subsection{Comparison to Planck 2018}

The Planck 2018 measured values \cite{Planck2020} are
$\Omega_\Lambda^{\rm Planck} = 0.6847 \pm 0.0073$ and
$\Omega_M^{\rm Planck} = 0.3153 \pm 0.0073$.
Our derived values match to:
\begin{equation*}
|\OmegaDE^{\rm derived} - \Omega_\Lambda^{\rm Planck}| =
|0.6843 - 0.6847| = 4 \times 10^{-4} = 0.06\%
\;\;(0.06\sigma\text{ relative to Planck error}).
\end{equation*}
The match is well within $1\sigma$ for both density fractions.
The derivation has no fitted parameters: the substrate
threshold $\Tstar = 5/7$, the wobble $W = 3/50$, and the
$\Mtwentytwo$-irreducible decomposition are all structural inputs from
the substrate algebra and representation theory, with no observational
data entering the derivation chain.

%%% ============================================================
\section{Relationship to the dynamical companion paper}\label{sec:companion}

% [TO FILL: explain that the companion paper takes Omega_DE as
% phenomenological input from Planck (= 0.6849); this paper supplies
% Omega_DE from substrate algebra; together they fix dark sector
% density (here) and dynamical evolution (companion).]

The companion paper \cite{SandersGishLogQuint2026} studies a
continuous scalar-field model with potential $V(\Xi) = \Lambda^4 \Xi
\log\Xi$ and analyzes its dynamical equation-of-state evolution
$w_{\rm DE}(z)$ over cosmic history, including a freeze-thaw transit
trajectory with Type-T, Type-F, and Type-A regimes. That paper takes
$\OmegaDE \approx 0.685$ as a phenomenological input from Planck 2018
constraints, with no microphysical derivation.

The present paper supplies the structural source for the same
$\OmegaDE = 479/700 = 0.6843$ value from substrate algebra alone,
with no observational input. The two papers are therefore
complementary: \cite{SandersGishLogQuint2026} provides the dynamical
content of the dark-energy sector (its $w(z)$ trajectory and
falsifiable Stage-IV signatures), while the present paper provides
the structural source of the asymptotic density fraction $\OmegaDE$.
Combined, the two papers fix both the density and the dynamics of
the dark-energy sector from a single underlying substrate algebra.

%%% ============================================================
\section{Falsifiable predictions}\label{sec:falsification}

% [TO FILL: state the bonus prediction tau_sigma ≈ 5.22 Gyr periodic
% modulation in cosmic structure formation rates; derive from sigma's
% 6-cycle projected to cosmological scale via the W/2 factor; identify
% Stage-IV survey tests (DESI DR3, Euclid, Rubin LSST) that could
% detect or rule out this signature.]

Beyond the $\OmegaDE$ value matching observation, the framework makes
a falsifiable structural prediction: a periodic modulation in cosmic
structure formation rates with period
\begin{equation}\label{eq:tau-sigma}
\tau_\sigma \;\approx\; 5.22 \,\textrm{Gyr},
\end{equation}
arising from the projection of the substrate's $\sigma$ $6$-cycle to
cosmological timescales via the $W/2$ factor (see derivation in
Appendix). This is a testable Stage-IV survey signature.

% [TO FILL: explicit derivation of tau_sigma from sigma cycle period
% and the W/2 cosmic-scale projection. Reference DESI DR3, Euclid,
% Rubin LSST sensitivity to oscillatory features in f sigma_8 and
% large-scale structure power spectra.]

%%% ============================================================
\section{Rigor demarcation}\label{sec:rigor}

We are explicit about the rigor level of each step in the derivation
chain.

\paragraph{Locked theorems (with proofs).}
\begin{enumerate}
\item Theorem \ref{thm:zeta-factorization}: $\zTIG(\Tstar)$
      pre-cancellation denominator equals $7^8 \cdot |\Mtwentytwo|$;
      $\Tstar$ is unique among $a/7$ fractions with this property.
\item The decomposition $\mathbb{C}^{22} \cong \Vtriv \oplus
      \Vtwentyone$ is a standard consequence of the double transitivity
      of the $\Mtwentytwo$ action on $\Omega_{22}$.
\item Proposition \ref{prop:equivariance}: $P_{\rm triv}$ and $P_{21}$
      commute with every element of $\Mtwentytwo$ (verified
      computationally over $1000$ random permutations and proven
      analytically).
\item The arithmetic identity $\OmegaDE = \Tstar - \Whalf = 479/700$
      and its match to Planck 2018 within $0.06\%$.
\end{enumerate}

\paragraph{Structurally motivated identifications (not yet theorems).}
\begin{enumerate}
\item The identification \eqref{eq:identification} of $G$ with $\Vtriv$
      and $K$ with $\Vtwentyone$ is justified by the dimensional identity
      $K = (1/7)\cdot (\dim \Vtwentyone /50)$ with all quantities entering
      from substrate algebra and representation theory. We do not have
      an independent first-principles derivation of why the wobble
      decomposes along these specific irreducibles; the identification
      is structurally consistent and quantitatively forced, but not
      yet a proven theorem from substrate dynamics alone.
\item The $50\%$ duty cycle of the wobble is established in
      \cite{SandersGishSigma} but the derivation of the duty cycle from
      first principles is itself structurally motivated rather than
      independently theoremized.
\end{enumerate}

\paragraph{Path to full theoremhood.}
% [TO FILL: outline the M_22-equivariant Hamiltonian formalism that
% would derive the V_trivial/V_21 split from substrate dynamics
% directly; outline the Steiner labeling derivation from substrate
% algebra alone (currently structurally motivated but technically a
% labeling choice). Both are concrete representation-theory
% follow-ups.]

%%% ============================================================
\section{Comparison to literature}\label{sec:literature}

% [TO FILL: brief comparison to (a) information-theoretic dark
% energy (Verlinde, Sheykhi); (b) sporadic-group physics applications
% (very rare; cite few examples like Conway-Norton moonshine,
% Mathieu-moonshine in K3 surfaces); (c) number-theoretic cosmology
% (rare); (d) standard quintessence (cite for completeness).]

%%% ============================================================
\section{Scope and limitations}\label{sec:scope}

% [TO FILL: explicit statement of what this paper claims and does
% not claim, parallel to \S 8 of the companion paper. Key claims:
% locked theorems above; key non-claims: full derivation of wobble
% structure from first principles, full theoremhood of the
% structural identification, etc.]

The present paper claims:
\begin{itemize}
\item The locked theorems and propositions enumerated in \S\ref{sec:rigor}.
\item The arithmetic identity $\OmegaDE = 479/700 = 0.6843$ matching
      Planck 2018 within $0.06\%$.
\item The structural identification \eqref{eq:identification} as
      consistent with all known substrate-algebraic and
      representation-theoretic data.
\end{itemize}

The present paper does not claim:
\begin{itemize}
\item A first-principles derivation of the wobble structure $W = 3/50$
      from substrate dynamics alone (this is established in
      \cite{SandersGishSigma} and \cite{SandersGishFourCore}).
\item A derivation of the $50\%$ duty cycle from first principles.
\item A full $\Mtwentytwo$-equivariant Hamiltonian description of
      the wobble dynamics.
\item Any derivation of the cosmological constant scale $\Lambda$ in
      physical units; only the dimensionless fraction $\OmegaDE$ is
      derived.
\end{itemize}

%%% ============================================================
\section*{Acknowledgments}

% [TO FILL]

\begin{thebibliography}{99}

\bibitem{Planck2020}
Planck Collaboration, \emph{Planck 2018 results. VI. Cosmological
parameters}, Astron. Astrophys. \textbf{641}, A6 (2020).

\bibitem{Weinberg1989}
S.~Weinberg, \emph{The cosmological constant problem}, Rev. Mod. Phys.
\textbf{61}, 1 (1989).

\bibitem{MartelShapiroWeinberg1998}
H.~Martel, P.~R.~Shapiro, and S.~Weinberg, \emph{Likely values of
the cosmological constant}, Astrophys. J. \textbf{492}, 29 (1998).

\bibitem{Verlinde2011}
E.~Verlinde, \emph{On the origin of gravity and the laws of Newton},
JHEP \textbf{04}, 029 (2011).

\bibitem{SandersGishLogQuint2026}
B.~R.~Sanders and M.~Gish, \emph{The Freeze-Thaw Transit: A Dual-Regime
Scalar Dark-Energy Model with an Analytic Turnaround at $e^{-1}$},
Companion paper, 2026.

\bibitem{SandersGishSigma}
B.~R.~Sanders and M.~Gish, \emph{Non-Associativity Decay in Binary
Composition Tables over $\mathbb{Z}/N\mathbb{Z}$}, Companion paper, 2026.

\bibitem{SandersGishFourCore}
B.~R.~Sanders and M.~Gish, \emph{Joint Closure, Per-Coordinate Fuse Data,
and a Closed-Form Algebraic Attractor of Two Commutative Binary
Operations on $\mathbb{Z}/10\mathbb{Z}$}, Companion paper, 2026.

\bibitem{ZenodoM22}
B.~R.~Sanders and M.~Gish, \emph{$\Mtwentytwo$ representation-theoretic
verification scripts and rigor documentation}, Zenodo deposit DOI~%
TBD (2026).

% [Additional references to fill: M_22 Atlas, Steiner system
% references, sporadic group monographs, Mathieu moonshine refs,
% Stage-IV survey papers]

\end{thebibliography}

\end{document}
